\section{Principle of the method}
Sometimes called \emph{Theory} or \emph{Theoretical background}.

Describe the theoretical principles the experiment is based on and how they are used to reach the objectives. Include every concept needed to understand the report. If the report belongs to a series, you may refer to earlier reports instead of repeating the same theory.

A glossary-style list can work well:
\begin{description}
    \item[Stoichiometry] The relationship between the relative quantities of substances taking part in a reaction or forming a compound, typically a ratio of whole numbers.
    \item[Atomic mass] The mass of an atom of a chemical element expressed in unified atomic mass units (\unit{\dalton}). It is approximately equal to the number of protons and neutrons in the atom (the mass number) or, for an element, to the average weighted by the natural abundance of its isotopes.
\end{description}

Alternatively, use \verb|\subsection{...}| to organise longer explanations, and number the key equations so that you can refer to them later with \verb|\cref{eq:label}|.
